% BEGIN MOTIVATION MEASUREMENT APPENDIX
\appendix
\section{Motivation Measurement Protocol}
\label{sec:motivation-protocol}

\paragraph{Scope and configuration.}
Figures~\ref{fig:motivation}(\subref{fig:motivation-recompute}) and
\ref{fig:motivation-recovery-exposure} use the same pinned
Qwen2.5-7B-Instruct-1M weights, SGLang 0.5.10rc0, PyTorch 2.9.1/CUDA 12.8,
TP2, BF16 weights/KV, and batch one on two A100 PCIe 40\,GB GPUs.
They use separate allocations of the same GPU type, with one fixed
repeated synthetic text tokenized by the model tokenizer. The traces
supply workflow positions and length proxies, not benchmark prompts.
Our adapter uses the pinned MInference vertical/slash configuration and
DCA position handling, with sparse attention above 32,768 tokens,
including the final one-token EXTEND. The SM80 sparse-kernel build and
model/attention sources are archived with hashes. GQA K/V are shared;
prefix execution groups at most two query heads, and a single query
batches all heads. These are instrumented worker measurements, rather
than throughput measurements of an unmodified serving frontend.

\paragraph{Common endpoint.}
For $C=M+1$, full prefill computes $M$ tokens from empty KV in chunks of
at most 8,192, then executes one token separately. Reuse consumes exactly
the same sparse-generated $M$-token state, including partial pages, and
executes that final token. Synchronized wall time is reduced by maximum
across the TP ranks and ends at model-body completion with all $C$ tokens'
KV resident, before the output vocabulary projection. Model loading,
backup creation, and correctness checks are outside the timer. Detection,
rerouting, serving-queue delay and an injected failure are absent.

\paragraph{Restore paths.}
HBM begins with resident KV. DRAM copies all historical layers from
pinned host buffers before computation. Shared storage opens, reads and
closes per-rank KV files on the compute node's DPC shared filesystem,
then performs the same H2D copy and computation. These stages do not
overlap; there is no compression or selective loading. The 1,010,000-token
case restores 57,917,382,656 logical history bytes across both ranks;
files additionally contain recorded alignment padding. Direct I/O and
page-residency checks exclude Linux file-page hits, but DPC server and
internal caches are uncontrolled. The curve is a shared-storage path
measurement, not cold-device latency or an identified SSD measurement.
KV files are removed after measurement.

\paragraph{Timing samples.}
Figure~\ref{fig:motivation}(\subref{fig:motivation-recompute}) uses one
allocation for 8,192, 65,536, 262,144 and 1,010,000 tokens. HBM/DRAM have
one warmup and three timings per length; full prefill and shared storage
have one timing each. Short forwards and operator checks precede timing;
64K also runs a complete original all-head sparse reference. There is
no extra complete-prefill warmup at 256K/1.01M or storage-read warmup.
Consequently, these full-prefill/storage points describe individual
runs, not stable latency distributions.
Figure~\ref{fig:motivation-recovery-exposure} uses two allocations at
$C=100{,}865$, $105{,}921$, $108{,}865$ and $113{,}665$, with one complete
warmup and three full-prefill timings at each length. The initial allocation completed three lengths; the last length was
resumed after an early termination signal. Its incomplete timings are
retained in the artifact but excluded as a whole. The figure reports
medians and min--max ranges. Repetitions are timing samples of the same input,
not independent workflows or confidence intervals. Original timings and
all samples are retained in the artifact.

\paragraph{Validation and workflow mapping.}
Complete-history hashes, final hidden states, all-layer new KV and greedy
outputs are checked across reconstruction/reuse paths. Repeated full
prefills additionally check output agreement, and CUDA-event timelines
check serialized copies. These checks validate matched sparse-state
execution, not equivalence to dense attention or long-context task quality.
For Figure~\ref{fig:motivation-recovery-exposure}, the four history lengths
are the active requests' trace hash-block counts multiplied by 64.
They are size proxies, not exact original-model token counts or observed
KV frontiers at failure. The example was selected for four parallel
branches and a gap near 100\,s; it is not its workflow's maximum-gap group.
Each counterfactual separately inserts a measured full-prefill duration
as an additional branch pause at 29.5\,s after the fork. The all-branch
join, other execution and queues stay fixed. Thus only reconstruction
cost is measured: the plotted WCT increment is a conditional calculation,
not observed fault recovery in the Claude Code workflow.
% END MOTIVATION MEASUREMENT APPENDIX
